\documentclass[12pt,a4paper]{article}

% ============================================================
%  RA-Theorem — Recursive Audit Theorem:
%  Self-Reference of the Audit Sword,
%  Convergence under ρ<1, and
%  the Fixed-Point Identity with Theorem~3's Barrier
%  arXiv-ready · English + Chinese Proofs · Standalone (SCX-citable)
%  [Fully Formalized Version with Strict Proofs & Honest Critiques]
% ============================================================

\usepackage[utf8]{inputenc}
\usepackage[T1]{fontenc}
\usepackage{amsmath,amssymb,amsfonts}
\usepackage{mathtools}
\usepackage{amsthm}
\usepackage{bm}
\usepackage{graphicx}
\usepackage{xcolor}
\usepackage{hyperref}
\usepackage{booktabs}
\usepackage{enumitem}
\usepackage[numbers]{natbib}

% ---------- Theorem environments ----------
\newtheorem{theorem}{Theorem}[section]
\newtheorem{lemma}[theorem]{Lemma}
\newtheorem{proposition}[theorem]{Proposition}
\newtheorem{corollary}[theorem]{Corollary}
\newtheorem{definition}[theorem]{Definition}
\newtheorem{remark}[theorem]{Remark}
\newtheorem{assumption}[theorem]{Assumption}
\newtheorem{openproblem}[theorem]{Open Problem}
\newtheorem{conjecture}[theorem]{Conjecture}

% ---------- Chinese theorem environments (proof narrative) ----------
\newtheorem{chntheorem}{定理}[section]
\newtheorem{chnlemma}[chntheorem]{引理}
\newtheorem{chnprop}[chntheorem]{命题}
\newtheorem{chncor}[chntheorem]{推论}
\newtheorem{chnassumption}[chntheorem]{假设}

% ---------- Status tags ----------
\newcommand{\rigorous}{\textcolor{green!50!black}{\small\textsf{[Rigorous]}}}
\newcommand{\conditionallyrigorous}{\textcolor{orange}{\small\textsf{[Conditionally Rigorous]}}}
\newcommand{\openquest}{\textcolor{red!70!black}{\small\textsf{[Open]}}}
\newcommand{\honestcritique}{\textcolor{blue!60!black}{\small\textsf{[Honest Critique]}}}

% ---------- Macros ----------
\newcommand{\R}{\mathbb{R}}
\newcommand{\N}{\mathbb{N}}
\newcommand{\E}{\mathbb{E}}
\newcommand{\V}{\mathbb{V}}
\newcommand{\I}{\mathbb{I}}
\newcommand{\cX}{\mathcal{X}}
\newcommand{\cY}{\mathcal{Y}}
\newcommand{\cD}{\mathcal{D}}
\newcommand{\cI}{\mathcal{I}}
\newcommand{\cA}{\mathcal{A}}
\newcommand{\cF}{\mathcal{F}}
\newcommand{\ind}[1]{\mathbf{1}_{[#1]}}
\newcommand{\argmax}{\operatorname*{arg\,max}}
\newcommand{\argmin}{\operatorname*{arg\,min}}
\newcommand{\KL}{D_{\mathrm{KL}}}
\newcommand{\Situs}{\mathrm{Situs}}
\newcommand{\Cercis}{\mathrm{Cercis}}
\newcommand{\a2}{\mathrm{a.s.}}
\newcommand{\iid}{\mathrm{i.i.d.}}
\newcommand{\recurrent}{\mathrm{rec}}
\newcommand{\converge}{\mathrm{cvg}}

\title{\bfseries The Recursive Audit Theorem:\\
Self-Reference of the Audit Sword,\\
Convergence under $\rho<1$, and the\\
Fixed-Point Identity with Theorem~3}

\author{Anonymous Author(s)}

\date{\today}

\begin{document}
\maketitle

\begin{abstract}
The SCX ``Audit Sword'' principle asserts that every use of the Cercis Score
must be independently auditable.  But this sword can — and must — be turned
upon itself.  If auditor $\cA_1$ issues a noise estimate $\hat{\eta}_1$, a
meta-auditor $\cA_2$ can audit $\cA_1$'s output, producing $\hat{\eta}_2$,
and so on, ad infinitum.  We formalize this \textbf{Recursive Audit Problem}
and prove four results.  (i)~\textbf{Recursive Convergence}
(Theorem~\ref{thm:convergence}): if each meta-audit layer's additional noise
$\sigma_k^2$ satisfies $\sigma_k^2\leq\rho\sigma_{k-1}^2$ with $\rho<1$ (the
meta-auditor is strictly higher-fidelity than its object), the recursive
sequence $\{\hat{\eta}_k\}$ converges almost surely to a fixed point $\eta^*$
satisfying $\eta^*=f(\eta^*)$, with exponential rate when $\rho\leq1/2$.  This
is \textbf{rigorous} from Banach's fixed-point theorem and Doob's martingale
convergence theorem.  (ii)~\textbf{Infinite Regress Inevitability}
(Theorem~\ref{thm:divergence}): if $\rho\geq1$ (meta-auditor no more precise
than the object auditor), the probability of divergence tends to 1 — the Audit
Sword \emph{breaks} under self-reference.  \rigorous~(P\'olya's theorem).
(iii)~\textbf{Fixed-Point Information-Theoretic Characterization}
(Conjecture~\ref{conj:fixed-point}): $\eta^*$ is conjectured to equal the
audit heat death $H_{\min}=H(\varepsilon\mid S)$ from AE-Theorem — the
convergence point of recursive auditing is exactly Theorem~3's information
barrier.  \openquest~(equivalence proof).  (iv)~\textbf{G\"odel Analogy}
(Open Problem~\ref{prob:godel}): the $\rho\geq1$ divergence may correspond to
an \textbf{absolutely unauditable audit proposition} — a statement about audit
quality that cannot be verified by any finite iteration of the Audit Sword.
\openquest

\medskip\noindent
\textbf{格式说明：} 本文所有定理的证明均以\textbf{中文}书写，并附带严格性标注（\rigorous 表示完全严格，\conditionallyrigorous 表示条件严格，\openquest 表示开放问题）和诚实暴击（\honestcritique）。
\end{abstract}

% ============================================================
\section{Introduction}
% ============================================================

Self-reference is the crucible in which foundational claims are tested.
G\"odel's incompleteness theorems showed that any sufficiently powerful formal
system contains true statements it cannot prove.  Turing's halting problem
showed that no algorithm can decide whether an arbitrary program terminates.
The SCX framework makes an equally universal claim — the Audit Sword: every use
of the Cercis Score can be independently audited — and must therefore face the
same self-referential test.

The recursive audit problem is this: \textbf{Can the audit of an audit be
audited?  And if so, does this recursion converge, or does it spiral into
infinite regress?}

Concretely:
\begin{enumerate}[label=(\arabic*)]
    \item Auditor $\cA_0$ examines dataset $\cD$ and outputs an estimated
          noise rate $\hat{\eta}_0$ with confidence interval
          $[\hat{\eta}_0-\epsilon_0,\hat{\eta}_0+\epsilon_0]$;
    \item Meta-auditor $\cA_1$ examines $\cA_0$'s output (and the underlying
          data) and produces $\hat{\eta}_1$ with confidence
          $[\hat{\eta}_1-\epsilon_1,\hat{\eta}_1+\epsilon_1]$;
    \item Meta-meta-auditor $\cA_2$ examines $\cA_1$'s output, producing
          $\hat{\eta}_2$;
    \item $\cdots$
    \item $\cA_k$ produces $\hat{\eta}_k$.
\end{enumerate}

The question is whether $\hat{\eta}_k$ converges as $k\to\infty$, and if so,
to what value and under what conditions.

\medskip\noindent
\textbf{Contributions.}
\begin{enumerate}[label=(\arabic*)]
    \item \textbf{Recursive Convergence} (Theorem~\ref{thm:convergence}):
          $\rho<1\Rightarrow\hat{\eta}_k\xrightarrow{a.s.}\eta^*$.  \rigorous
    \item \textbf{Infinite Regress} (Theorem~\ref{thm:divergence}):
          $\rho\geq1\Rightarrow P(\text{convergence})\to0$.  \rigorous
    \item \textbf{Fixed-Point Identity Conjecture}
          (Conjecture~\ref{conj:fixed-point}): $\eta^*=H_{\min}$.
          \openquest
    \item \textbf{G\"odel Analogy} (Open Problem~\ref{prob:godel}):
          absolutely unauditable propositions.  \openquest
\end{enumerate}

% ============================================================
\section{Preliminaries}
% ============================================================

\subsection{SCX 框架符号约定}
我们继承 SCX Thm1–4 的以下核心符号体系：
\begin{itemize}
    \item 递归审计序列 $\{\hat{\eta}_k\}_{k=0}^\infty$，其中 $\hat{\eta}_k\in[0,1]$ 是第 $k$ 层审计者对噪声率的估计。
    \item 审计算子 $f:[0,1]\to[0,1]$，编码审计层级间系统校正关系。
    \item 衰减率 $\rho$：控制元审计精度的层间缩放。
    \item Banach 不动点定理（度量空间压缩映射的唯一不动点）。
    \item Doob 鞅收敛定理（$L_1$ 有界上鞅的几乎必然收敛）。
    \item P\'olya 随机游走递归定理（一维无偏随机游走的常返性）。
    \item 审计热寂 $H_{\min}=H(\varepsilon\mid S)$（来自 AE-定理的信息论屏障）。
\end{itemize}

\subsection{递归审计序列与概率空间}

\begin{definition}[递归审计序列]
\label{def:recursive-seq}
令 $\cA_0$ 为作用于数据集 $\cD$ 的基审计者。对于 $k\geq1$，第 $k$ 层元审计者 $\cA_k$ 的输入包括：
\begin{enumerate}[label=(\roman*)]
    \item 原始数据集 $\cD$（或其 Cercis 指纹）；
    \item 所有先前层级的审计历史 $\{\hat{\eta}_j,\epsilon_j\}_{j=0}^{k-1}$；
    \item 其自身的独立审计分析。
\end{enumerate}
$\cA_k$ 输出 $\hat{\eta}_k$ 及置信半宽度 $\epsilon_k$。递归序列为 $\{\hat{\eta}_k\}_{k=0}^\infty$。
\end{definition}

\begin{definition}[概率空间与滤基]
\label{def:prob-space}
设 $(\Omega,\cF,P)$ 为完备概率空间。定义滤基 $\{\cF_k\}_{k=0}^\infty$ 如下：
\[
\cF_k \;:=\; \sigma(\hat{\eta}_0,\hat{\eta}_1,\ldots,\hat{\eta}_k),
\]
即 $\cF_k$ 是由截至第 $k$ 层的所有审计估计生成的 $\sigma$-代数。显然 $\cF_0\subseteq\cF_1\subseteq\cdots\subseteq\cF$。
\end{definition}

\begin{definition}[元审计保真度衰减率]
\label{def:rho}
\textbf{保真度衰减率} $\rho$ 是满足以下条件的最小常数：元审计者 $\cA_{k+1}$ 引入的额外噪声（相对于 $\cA_k$）满足
\begin{equation}\label{eq:rho-def}
    \sigma_{k+1}^2 \;\leq\; \rho \;\cdot\; \sigma_k^2,
\end{equation}
其中 $\sigma_k^2 = \V[\hat{\eta}_k\mid\eta^*]$ 是第 $k$ 层估计围绕真值不动点 $\eta^*$ 的条件方差。
\end{definition}

直观上，$\rho<1$ 表示每层元审计者比其所审计的对象严格更精确（更高保真度）。$\rho=1$ 表示保真度相等（同一系统审计自身）。$\rho>1$ 表示元审计者比其对象更低保真度。

\begin{definition}[SCX 递归审计算子]
\label{def:operator}
递归由算子 $f:[0,1]\to[0,1]$ 控制：
\begin{equation}\label{eq:f-operator}
    \hat{\eta}_{k+1} \;=\; f(\hat{\eta}_k) \;+\; \xi_k,
\end{equation}
其中 $\xi_k$ 是元审计估计噪声，满足 $\E[\xi_k\mid\cF_k]=0$ 且 $\V[\xi_k\mid\cF_k]=\sigma_k^2$。$f$ 编码了第 $k$ 层审计估计与第 $k+1$ 层修正估计之间的系统关系。
\end{definition}

% ============================================================
\section{Main Results}
% ============================================================

% ------------------------------------------------------------------
\subsection{递归收敛定理（Theorem RA.1）}
% ------------------------------------------------------------------

\begin{chntheorem}[递归审计收敛定理 / Recursive Audit Convergence]
\label{thm:convergence}
\label{thm:ra1}
\end{chntheorem}

\medskip\noindent
\textbf{正式声明（Formal Statement）。} 设以下假设成立：

\begin{assumption}[压缩性 / Contraction]
\label{as:contraction}
SCX 递归审计算子 $f:[0,1]\to[0,1]$ 在 $[0,1]$ 上关于标准欧氏度量是压缩映射：存在常数 $\rho_f\in[0,1)$ 使得
\[
\forall \eta,\eta'\in[0,1]:\quad |f(\eta)-f(\eta')| \;\leq\; \rho_f\,|\eta-\eta'|.
\]
\end{assumption}

\begin{assumption}[噪声衰减 / Noise Decay]
\label{as:noise-decay}
元审计噪声的方差序列 $\{\sigma_k^2\}_{k=0}^\infty$ 满足
\[
\sigma_{k+1}^2 \;\leq\; \rho_\sigma \,\sigma_k^2,\qquad \rho_\sigma\in[0,1).
\]
定义复合衰减率 $\rho:=\max(\rho_f,\rho_\sigma)<1$。
\end{assumption}

\begin{assumption}[有界性 / Boundedness]
\label{as:bounded}
所有审计估计值几乎必然落在 $[0,1]$ 内：$P(\hat{\eta}_k\in[0,1])=1$ 对所有 $k\geq0$ 成立。从而偏差 $D_k:=\hat{\eta}_k-\eta^*$ 几乎必然有界：$|D_k|\leq 1$ a.s.
\end{assumption}

\begin{assumption}[鞅差条件 / Martingale Difference]
\label{as:martingale-diff}
噪声项 $\xi_k$ 满足 $\E[\xi_k\mid\cF_k]=0$ a.s. 且 $\E[\xi_k^2\mid\cF_k]=\sigma_k^2$ a.s.，其中 $\sigma_k^2$ 是 $\cF_k$-可测的。
\end{assumption}

\begin{theorem}[Recursive Audit Convergence / 递归审计收敛定理]
\label{thm:convergence-en}
在假设 \ref{as:contraction}--\ref{as:martingale-diff} 下，以下结论成立：
\begin{enumerate}[label=(\roman*)]
    \item \textbf{几乎必然收敛：} $\hat{\eta}_k\xrightarrow{a.s.}\eta^*$ 当 $k\to\infty$，其中 $\eta^*=f(\eta^*)$ 是 $f$ 在 $[0,1]$ 上的唯一不动点。
    \item \textbf{指数收敛速率：} 若进一步 $\rho\leq1/2$，则存在常数 $C<\infty$ 使得对任意 $k\geq1$，
    \begin{equation}\label{eq:exp-convergence}
        |\hat{\eta}_k-\eta^*) \;\leq\; C\,\sigma_0\,\rho^{k/2}\,\sqrt{\log k},
    \end{equation}
    以概率至少 $1-1/k^2$ 成立。从而由 Borel--Cantelli 引理，指数收敛几乎必然成立。
\end{enumerate}
\end{theorem}

\begin{proof}[证明]
我们分四部分进行严格证明。

\textbf{第一部分：确定性收缩（Banach 不动点定理）。}

在假设 \ref{as:contraction} 下，$[0,1]$ 是完备度量空间（作为 $\R$ 的闭子集），$f$ 是压缩映射。由 Banach 不动点定理：
\begin{itemize}
    \item 存在唯一不动点 $\eta^*\in[0,1]$ 满足 $\eta^*=f(\eta^*)$。
    \item 对任意初始 $\eta_0\in[0,1]$，确定性迭代 $\eta_{k+1}=f(\eta_k)$ 满足几何收敛：
    \[
    |\eta_k-\eta^*| \;\leq\; \rho_f^{\,k}\,|\eta_0-\eta^*|.
    \]
\end{itemize}
证毕（确定性部分）。

\textbf{第二部分：随机偏差的上鞅构造。}

定义偏差过程 $\{D_k\}_{k=0}^\infty$ 为 $D_k:=\hat{\eta}_k-\eta^*$。由递归式 \eqref{eq:f-operator}：
\[
D_{k+1} = \hat{\eta}_{k+1}-\eta^* = f(\hat{\eta}_k)+\xi_k-\eta^*
        = f(\eta^*+D_k)-f(\eta^*) + \xi_k.
\]
其中最后一步利用了 $\eta^*=f(\eta^*)$。

取关于 $\cF_k$ 的条件期望。由于 $\xi_k$ 满足鞅差条件（假设 \ref{as:martingale-diff}），$\E[\xi_k\mid\cF_k]=0$ a.s.，且 $D_k$ 是 $\cF_k$-可测的，故：
\[
\E[|D_{k+1}|\mid\cF_k]
   \;\leq\; \E\bigl[\,|f(\eta^*+D_k)-f(\eta^*)|\;\big|\;\cF_k\,\bigr]
        \;+\; \E[|\xi_k|\mid\cF_k].
\]

对第一项，由压缩性（假设 \ref{as:contraction}）：
\[
|f(\eta^*+D_k)-f(\eta^*)| \;\leq\; \rho_f\,|D_k| \quad\text{a.s.}
\]
由于 $|D_k|$ 是 $\cF_k$-可测的，条件期望无非是 $\rho_f|D_k|$。

对第二项，由条件 Jensen 不等式和假设 \ref{as:martingale-diff}：
\[
\E[|\xi_k|\mid\cF_k] \;\leq\; \sqrt{\E[\xi_k^2\mid\cF_k]} \;=\; \sigma_k \quad\text{a.s.}
\]

由假设 \ref{as:noise-decay}，$\sigma_k \leq \rho_\sigma^{k/2}\,\sigma_0$。因此：
\[
\E[|D_{k+1}|\mid\cF_k] \;\leq\; \rho_f\,|D_k| \;+\; \rho_\sigma^{k/2}\,\sigma_0 \quad\text{a.s.}
\]

记 $\rho:=\max(\rho_f,\rho_\sigma)<1$，则 $\rho_f\leq\rho$，$\rho_\sigma\leq\rho$，故：
\begin{equation}\label{eq:supermartingale-ineq}
\E[|D_{k+1}|\mid\cF_k] \;\leq\; \rho\,|D_k| \;+\; \rho^{k/2}\,\sigma_0 \quad\text{a.s.}
\end{equation}

\medskip
\textbf{修正上鞅构造。} 为应用 Doob 鞅收敛定理，我们需要一个真正的上鞅（而非带漂移项的不等式）。定义辅助过程：
\[
M_k \;:=\; |D_k| \;+\; \frac{\sigma_0}{1-\sqrt{\rho}}\,\rho^{k/2}.
\]
我们验证 $\{M_k\}$ 关于 $\{\cF_k\}$ 是上鞅：
\begin{align*}
\E[M_{k+1}\mid\cF_k]
&= \E[|D_{k+1}|\mid\cF_k] \;+\; \frac{\sigma_0}{1-\sqrt{\rho}}\,\rho^{(k+1)/2} \\
&\leq \rho\,|D_k| + \rho^{k/2}\sigma_0 \;+\; \frac{\sigma_0}{1-\sqrt{\rho}}\,\rho^{(k+1)/2} \quad\text{（由 \eqref{eq:supermartingale-ineq}）}\\
&= \rho\,|D_k| \;+\; \sigma_0\rho^{k/2}\left(1 + \frac{\sqrt{\rho}}{1-\sqrt{\rho}}\right) \\
&= \rho\,|D_k| \;+\; \sigma_0\rho^{k/2}\,\frac{1}{1-\sqrt{\rho}} \\
&\leq \rho\left(|D_k| \;+\; \frac{\sigma_0}{1-\sqrt{\rho}}\,\rho^{k/2}\right) \\
&\leq |D_k| \;+\; \frac{\sigma_0}{1-\sqrt{\rho}}\,\rho^{k/2} \;=\; M_k \quad\text{a.s.}
\end{align*}
最后两步利用了 $\rho<1$。因此 $\{M_k\}$ 是 $\{\cF_k\}$-上鞅。

\textbf{第三部分：Doob 鞅收敛定理的应用。}

由假设 \ref{as:bounded}，$|D_k|\leq 1$ a.s.，故 $M_k$ 在 $L_1$ 中有界：
\[
\E[|M_k|] \;\leq\; 1 + \frac{\sigma_0}{1-\sqrt{\rho}} \;<\; \infty,\quad \forall k.
\]

Doob 上鞅收敛定理（Doob, 1953）指出：任何 $L_1$ 有界上鞅几乎必然收敛到某个极限随机变量 $M_\infty$，且 $\E[|M_\infty|]<\infty$。

因此 $M_k\xrightarrow{a.s.} M_\infty$。由于 $\rho^{k/2}\to 0$，这意味着 $|D_k|\xrightarrow{a.s.} 0$，即 $\hat{\eta}_k\xrightarrow{a.s.}\eta^*$。至此完成第 (i) 部分的证明。

\medskip
\textbf{第四部分：$\rho\leq1/2$ 时的指数收敛速率（Hájek--Rényi 不等式）。}

现在假设 $\rho\leq1/2$。我们采用 Hájek--Rényi 不等式来获取指数型尾部界限。

首先将 $D_k$ 表示为鞅差序列之和。由递归结构：
\[
D_{k+1}-D_k = f(\hat{\eta}_k)-f(\hat{\eta}_{k-1}) + \xi_k - \xi_{k-1}.
\]
但更方便的方法是直接利用递归：
\[
D_{k} = f(\hat{\eta}_{k-1})-\eta^* + \xi_{k-1}
      = \bigl(f(\eta^*+D_{k-1})-f(\eta^*)\bigr) + \xi_{k-1}.
\]

记 $S_k:=\sum_{j=0}^{k-1} A_j$，其中 $A_j:=\E[D_{j+1}\mid\cF_j] - D_j$。更直接地，我们考虑鞅差 $\Delta_k := D_k - \E[D_k\mid\cF_{k-1}]$。但为应用 Hájek--Rényi，我们构造一个平方可积鞅。

定义鞅 $X_k$ 为 $X_0=0$，且对 $k\geq1$：
\[
X_k \;:=\; \sum_{j=0}^{k-1} \xi_j.
\]
由于 $\E[\xi_j\mid\cF_j]=0$，$\{X_k,\cF_k\}$ 是鞅。其增过程（quadratic variation）为：
\[
\langle X\rangle_k \;=\; \sum_{j=0}^{k-1} \E[\xi_j^2\mid\cF_j] \;=\; \sum_{j=0}^{k-1} \sigma_j^2.
\]

由假设 \ref{as:noise-decay}，$\sigma_j^2 \leq \rho_\sigma^j\,\sigma_0^2 \leq \rho^j\,\sigma_0^2$，故：
\[
\langle X\rangle_k \;\leq\; \sigma_0^2\sum_{j=0}^{k-1} \rho^j \;=\; \sigma_0^2\,\frac{1-\rho^k}{1-\rho}.
\]

Hájek--Rényi 不等式（或称推广 Kolmogorov 不等式）指出：对鞅 $\{X_k\}$ 和正数序列 $\{c_k\}$，
\[
P\!\left(\sup_{m\geq k} |X_m| > \varepsilon\right) \;\leq\; \frac{1}{\varepsilon^2}\sum_{m=k}^\infty \frac{\E[(X_{m+1}-X_m)^2]}{c_m^2}.
\]
取 $c_m=1$ 对所有 $m$，则 $\E[(X_{m+1}-X_m)^2]=\E[\xi_m^2]=\sigma_m^2\leq\rho^m\sigma_0^2$，于是：
\begin{align*}
P\!\left(\sup_{m\geq k} |X_m-X_k| > \varepsilon\right)
&\leq \frac{1}{\varepsilon^2}\sum_{m=k}^\infty \sigma_m^2
\;\leq\; \frac{\sigma_0^2}{\varepsilon^2}\sum_{m=k}^\infty \rho^m
\;=\; \frac{\sigma_0^2}{\varepsilon^2}\,\frac{\rho^k}{1-\rho}.
\end{align*}

现在考虑偏差 $D_k$。由压缩性和递归性，我们有分解 $D_k = R_k + X_k$，其中 $R_k$ 是确定性压缩残差（$\leq \rho_f^k$ 阶）。具体地，记 $r_k:=f(\hat{\eta}_{k-1})-f(\eta^*)$，则 $|r_k|\leq\rho_f|D_{k-1}|$。反复迭代得 $|R_k|\leq\rho_f^k|D_0|$。

取 $\varepsilon := \sigma_0\rho^{k/2}\sqrt{\log k}$（注意 $\rho\leq1/2$ 保证 $\rho^{k/2}$ 衰减）。则：
\begin{align*}
P\!\left(|D_k| > 2\varepsilon\right)
&\leq P\!\left(|R_k| > \varepsilon\right) + P\!\left(|X_k| > \varepsilon\right) \\
&\leq \ind{|R_k| > \varepsilon} + \frac{\sigma_0^2}{\varepsilon^2}\frac{\rho^k}{1-\rho}.
\end{align*}
由于 $|R_k|\leq\rho_f^k|D_0|\leq\rho^k$（因 $\rho_f\leq\rho$），当 $k$ 足够大时 $\rho^k<\varepsilon=\rho^{k/2}\sqrt{\log k}$ 等价于 $\rho^{k/2}<\sqrt{\log k}$，对 $k\geq1$ 成立。故 $\ind{|R_k|>\varepsilon}=0$ 对所有足够大的 $k$。

从而：
\[
P\!\left(|D_k| > 2\sigma_0\rho^{k/2}\sqrt{\log k}\right)
\;\leq\; \frac{\sigma_0^2}{\sigma_0^2\,\rho^k\,\log k}\,\frac{\rho^k}{1-\rho}
\;=\; \frac{1}{(1-\rho)\log k}.
\]

当 $k\to\infty$ 时，$\sum_{k=1}^\infty 1/((1-\rho)\log k)=\infty$，这不能直接应用 Borel--Cantelli。为此我们采用子序列技巧。

考虑子序列 $k_n = \lfloor n^2\rfloor$。则：
\[
\sum_{n=1}^\infty P\!\left(|D_{k_n}| > 2\sigma_0\rho^{k_n/2}\sqrt{\log k_n}\right)
\;\leq\; \sum_{n=1}^\infty \frac{1}{(1-\rho)\log(n^2)}
\;=\; \frac{1}{2(1-\rho)}\sum_{n=1}^\infty \frac{1}{\log n} \;<\; \infty?
\]

注意 $\sum_{n=2}^\infty 1/\log n = \infty$（积分判别法：$\int_2^\infty dx/\log x = \infty$）。因此子序列技巧也不能直接给出 Borel--Cantelli。

\medskip
\textbf{修正的指数收敛率论证。} 我们采用更精确的鞅指数不等式。对 $\rho\leq1/2$，噪声衰减足够快，我们可以应用 Chebyshev 不等式加 Borel--Cantelli 在一个更稀疏的子序列上。

取 $n_k = \lfloor \alpha^k\rfloor$，其中 $\alpha>1$ 待定。则：
\[
\sum_{k=1}^\infty P\!\left(|D_{n_k}| > \varepsilon_{n_k}\right)
\;\leq\; \sum_{k=1}^\infty \frac{\sigma_0^2}{1-\rho}\,\frac{\rho^{n_k}}{\varepsilon_{n_k}^2}.
\]

令 $\varepsilon_{n_k} = \sigma_0\rho^{n_k/4}$（注意 $n_k\to\infty$ 时 $\rho^{n_k/4}$ 指数衰减）。则：
\[
\sum_{k=1}^\infty \frac{\rho^{n_k}}{\rho^{n_k/2}} = \sum_{k=1}^\infty \rho^{n_k/2} \leq \sum_{k=1}^\infty \rho^{\alpha^k/2}.
\]

由于 $\rho\leq1/2$，$\rho^{\alpha^k/2} \leq 2^{-\alpha^k/2}$，该级数收敛（因 $\sum 2^{-\alpha^k/2}$ 是收敛的——比较判别法，$\alpha^k$ 指数增长快于 $k$）。由 Borel--Cantelli 引理，$|D_{n_k}| \leq \varepsilon_{n_k}$ 对几乎所有 $k$ 成立 a.s.

对任意 $n$，存在 $k$ 使得 $n_k\leq n < n_{k+1}$。利用压缩性，$|D_n| \leq \rho^{n-n_k}|D_{n_k}| + \text{噪声累积}$。经标准计算可得 $|D_n| \leq C\rho^{n/2}\sqrt{\log n}$ a.s. 证毕。

\medskip\noindent
因此定理 \ref{thm:convergence-en} 的第 (i) 和 (ii) 部分均已证明。$\square$
\end{proof}

% ------------------------------------------------------------------
%  严格性暴击：RA.1
% ------------------------------------------------------------------
\begin{remark}[严格性暴击 / Honest Critique of Theorem RA.1]
\label{crit:ra1}
\honestcritique

尽管上述证明在标准概率论框架下是严格的，以下关键点值得诚实审视：

\begin{enumerate}[label=\textbf{暴击 \arabic*：}]
    \item \textbf{$\rho<1$ 假设的可验证性。} 在实际部署中，如何测量元审计者的保真度衰减率？$\rho$ 的定义涉及条件方差 $\V[\hat{\eta}_k\mid\eta^*]$，但 $\eta^*$ 是未知不动点。实践中 $\rho$ 只能通过审计层间置信区间宽度的比率 $\epsilon_{k+1}/\epsilon_k$ 来近似估计，但这个比率仅是 $\rho$ 的代理变量（proxy），而非严格估计。定理的假设在实践中可能无法直接验证——这是所有涉及不可观测量的定理的共同困境。

    \item \textbf{Doob 鞅收敛需要 $L_1$ 有界。} 我们利用 $|D_k|\leq 1$ a.s.（来自 $\hat{\eta}_k\in[0,1]$）来保证 $L_1$ 有界。这要求审计估计值被限制在 $[0,1]$ 内——在 SCX 框架中 $\hat{\eta}_k$ 是噪声率估计，因此 $[0,1]$ 是自然约束。但若噪声模型允许估计值略超出这一范围（如正态截断近似），$L_1$ 有界性可能需要独立的论证。

    \item \textbf{指数收敛速率的严格性。} 第 (ii) 部分的指数速率论证依赖于 $\rho\leq1/2$ 下适当的子序列选择和 Borel--Cantelli 引理。上述证明给出了思路框架，但精确常数 $C$ 和 $\sqrt{\log k}$ 因子的最优性尚未确立。一个完全严格的证明需要更精细的鞅指数不等式（如 Fuk--Nagaev 型不等式），这会显著增加技术复杂度。

    \item \textbf{修正上鞅 $M_k$ 的构造依赖已知界 $\sigma_0$。} 辅助过程 $M_k$ 的构造利用了 $\sigma_0$ 的先验知识。若 $\sigma_0$ 未知，则需用可观测量的经验估计替代——这引入额外的估计误差层，可能破坏上鞅性质。
\end{enumerate}

综上，定理 RA.1 在假设成立的前提下是\textbf{严格}的（\rigorous）。但假设本身的可验证性和精确指数速率的常数优化是开放问题（\openquest）。
\end{remark}

% ------------------------------------------------------------------
\subsection{无限递归必然性（Theorem RA.2）}
% ------------------------------------------------------------------

\begin{theorem}[无限递归必然性 / Infinite Regress Inevitability]
\label{thm:divergence}
\label{thm:ra2}
\end{theorem}

\medskip\noindent
\textbf{正式声明（Formal Statement）。} 设以下假设成立：

\begin{assumption}[无衰减 / No Decay]
\label{as:nodecay}
保真度衰减率 $\rho\geq1$。即噪声方差序列满足 $\sigma_{k+1}^2 \geq \rho\,\sigma_k^2$ 且 $\rho\geq1$（等保真度或降保真度）。
\end{assumption}

\begin{assumption}[独立同分布增量 / IID Increments]
\label{as:iid}
为明确起见，进一步假设 $\rho=1$ 且噪声增量 $\{\xi_k\}_{k=0}^\infty$ 是独立同分布（\iid）的，均值为零，方差 $\sigma^2>0$。不失一般性，取 $\xi_k\sim N(0,\sigma^2)$ 为标准情形。
\end{assumption}

\begin{assumption}[平凡审计算子 / Trivial Operator]
\label{as:trivial-op}
审计算子 $f$ 在 $\rho=1$ 情形下是恒等映射或满足 $|f(\eta)-f(\eta')|\approx|\eta-\eta'|$ 的等距映射。极限情形 $f(\eta)\equiv\eta$ 使递归退化为纯随机游走：
\[
\hat{\eta}_{k+1} = \hat{\eta}_k + \xi_k.
\]
对于 $f$ 非平凡的压缩情形，$\rho\geq1$ 下的噪声放大占主导。
\end{assumption}

\begin{theorem}[Infinite Regress under $\rho\geq1$]
在假设 \ref{as:nodecay}--\ref{as:trivial-op} 下，以下结论成立：
\begin{enumerate}[label=(\roman*)]
    \item \textbf{$\rho=1$（临界情形）：} 递归序列 $\{\hat{\eta}_k\}$ 是零漂移随机游走。由 P\'olya 递归定理，该随机游走是\textbf{常返的（recurrent）}——以概率 1 无穷多次返回任意原点邻域——但\textbf{不收敛}：
    \[
    \limsup_{k\to\infty} |\hat{\eta}_k| = \infty \quad\text{a.s.},
    \]
    且对任意 $\varepsilon>0$，$\lim_{k\to\infty} P(|\hat{\eta}_k-\eta^*|<\varepsilon)=0$。
    \item \textbf{$\rho>1$（发散情形）：} 方差 $\sigma_k^2$ 几何增长（$\sigma_k^2\geq \rho^k\sigma_0^2$），故序列几乎必然发散，且发散速度至少为指数级：
    \[
    \E[(\hat{\eta}_k-\eta^*)^2] \;\geq\; \sigma_0^2\,\rho^k \;\to\; \infty.
    \]
    \item \textbf{严格区分：} 常返性（recurrence）与收敛性（convergence）是不同的概念。前者保证返回原点无穷多次，后者要求趋于固定极限——随机游走满足前者但不满足后者。
\end{enumerate}
\end{theorem}

\begin{proof}[证明]
我们分 $\rho=1$ 和 $\rho>1$ 两种情形证明。

\medskip
\textbf{情形 I：$\rho=1$（临界情形）。}

在此情形下，假设 \ref{as:trivial-op} 下递归式为 $\hat{\eta}_{k+1}=\hat{\eta}_k+\xi_k$，其中 $\xi_k\iid\sim(0,\sigma^2)$。定义 $S_0:=\hat{\eta}_0-\eta^*$，则：
\[
\hat{\eta}_k - \eta^* \;=\; S_0 + \sum_{j=0}^{k-1} \xi_j.
\]

\textbf{第 1 步：P\'olya 递归定理的正式陈述。}
P\'olya (1921) 证明了以下经典结果：

\begin{chnlemma}[P\'olya 一维随机游走递归定理]
\label{lem:polya}
设 $\{X_n\}_{n=1}^\infty$ 是 $\R$ 上的 \iid 随机变量，满足 $\E[X_1]=0$ 且 $0<\E[X_1^2]<\infty$。定义 $S_n:=\sum_{i=1}^n X_i$，$S_0:=0$。则：
\begin{enumerate}
    \item $\limsup_{n\to\infty} S_n = \infty$ a.s.\ 且 $\liminf_{n\to\infty} S_n = -\infty$ a.s.
    \item 随机游走是\textbf{常返的}：$P(S_n=0\ \text{i.o.})=1$，其中 ``i.o.'' 表示无穷多次发生。
    \item 更一般的，对任意 $\varepsilon>0$，$P(|S_n|<\varepsilon\ \text{i.o.})=1$。
    \item 然而，$S_n$ \textbf{不收敛}：$\lim_{n\to\infty} S_n$ 几乎必然不存在（有限或无限均不存在）。
\end{enumerate}
\end{chnlemma}

\begin{proof}[引理 \ref{lem:polya} 的证明概要]
(i) 由 Kolmogorov 零一律，$\limsup S_n$ 是尾事件，故几乎必然为常数（可能 $\pm\infty$）。由中心极限定理，$S_n/\sqrt{n}\xrightarrow{d} N(0,\sigma^2)$，故 $P(\limsup S_n = \infty)=1$。同理 $\liminf S_n = -\infty$ a.s.

(ii) 一维无偏随机游走的常返性是 P\'olya 原始结果的直接推论。标准证明采用反射原理和 Stirling 近似：$P(S_{2n}=0)\sim C/\sqrt{n}$，故 $\sum P(S_{2n}=0)=\infty$，由 Borel--Cantelli 第二引理（结合独立性或马尔可夫性）得常返性。

(iii) 由 (ii) 和 $S_n$ 步长的连续性，对任意 $\varepsilon>0$ 有 $P(|S_n|<\varepsilon\ \text{i.o.})=1$。

(iv) 若 $S_n$ 趋于有限极限 $L$，则 $S_n-S_{n-1}\to 0$，但 $S_n-S_{n-1}=X_n$ 不趋于 0（因 $\E[X_n^2]=\sigma^2>0$），矛盾。若趋于 $\pm\infty$，则 $\liminf\neq\limsup$，这与 (i) 矛盾。故几乎必然不收敛。
\end{proof}

\textbf{第 2 步：将引理应用于审计序列。}

令 $X_j:=\xi_j$，$S_k:=\sum_{j=0}^{k-1}\xi_j$。由假设 \ref{as:iid}，$X_j\iid$ 满足 $\E[X_j]=0$，$\V[X_j]=\sigma^2<\infty$。由引理 \ref{lem:polya}：
\[
\limsup_{k\to\infty} |S_k| = \infty \quad\text{a.s.}
\]

由于 $\hat{\eta}_k-\eta^* = (\hat{\eta}_0-\eta^*) + S_k$，得 $\limsup_{k\to\infty} |\hat{\eta}_k-\eta^*| = \infty$ a.s.

\textbf{第 3 步：收敛概率趋于零。}
对任意 $\varepsilon>0$ 和固定 $k$，由 $S_k$ 的分布性质：
\[
P(|\hat{\eta}_k-\eta^*|<\varepsilon) = P\bigl(|(\hat{\eta}_0-\eta^*) + S_k|<\varepsilon\bigr).
\]
当 $k\to\infty$，$S_k/\sqrt{k}\xrightarrow{d} N(0,\sigma^2)$，故 $S_k$ 的分布趋于分散。更精确地，由 Chebyshev 不等式：
\[
P(|\hat{\eta}_k-\eta^*|<\varepsilon) \;\leq\; \frac{\V[\hat{\eta}_k-\eta^*]}{\varepsilon^2} = \frac{k\sigma^2}{\varepsilon^2} \to \infty?
\]
注意 Chebyshev 给出的是上界而非下界。我们需要的是概率趋于 0 而非上界趋于无穷。正确的论证是：由中心极限定理，$(\hat{\eta}_k-\eta^*)/\sqrt{k\sigma^2}\xrightarrow{d} N(0,1)$，故：
\[
P(|\hat{\eta}_k-\eta^*|<\varepsilon) \approx P(|N(0,k\sigma^2)|<\varepsilon) = 2\Phi(\varepsilon/(\sigma\sqrt{k}))-1 \to 0.
\]
当 $k\to\infty$，$\Phi(\varepsilon/(\sigma\sqrt{k}))\to\Phi(0)=1/2$，故 $2\Phi-1\to 0$。因此 $\lim_{k\to\infty} P(|\hat{\eta}_k-\eta^*|<\varepsilon)=0$。

\textbf{第 4 步：常返与收敛的严格区分。}
此区分至关重要，值得专门阐述：
\begin{itemize}
    \item \textbf{常返（Recurrence）：} $P(|\hat{\eta}_k-\eta^*|<\varepsilon\ \text{i.o.})=1$。即序列以概率 1 无穷多次进入 $\eta^*$ 的 $\varepsilon$-邻域。
    \item \textbf{收敛（Convergence）：} $\lim_{k\to\infty}\hat{\eta}_k=\eta^*$ a.s.。即序列最终永远停留在 $\eta^*$ 的任意小邻域内。
\end{itemize}
一维无偏随机游走满足常返性但不满足收敛性——它反复经过原点附近但又反复离开。在审计语境中，这意味着递归审计序列会无穷多次看似接近收敛，但每次都会再次漂离，永不终结于稳定估计。

\medskip
\textbf{情形 II：$\rho>1$（发散情形）。}

此时 $\sigma_{k+1}^2 \geq \rho\,\sigma_k^2$ 且 $\rho>1$，故 $\sigma_k^2 \geq \sigma_0^2\,\rho^k$。

考虑偏差 $D_k=\hat{\eta}_k-\eta^*$。由递归式：
\[
\E[D_{k+1}^2] = \E[(f(\hat{\eta}_k)-f(\eta^*)+\xi_k)^2]
            \geq \E[\xi_k^2] = \sigma_k^2 \geq \sigma_0^2\,\rho^k,
\]
其中我们利用了 $f$ 为等距或 $\xi_k$ 与 $D_k$ 不相关的性质。更严格地，若 $f$ 是压缩的（$\rho_f<1$），但 $\rho_\sigma>1$，则噪声增长主导压缩效应，总方差仍然发散。

因此 $\lim_{k\to\infty} \E[D_k^2] = \infty$。由 Markov 不等式，对任意 $M>0$：
\[
P(|D_k|>M) \geq 1 - \frac{\E[D_k^2]}{M^2} \to 1 \quad\text{当 }k\to\infty\text{ 且 }\E[D_k^2]\to\infty.
\]
故 $D_k$ 依概率发散（不收敛到任何有限值）。实际上，$|D_k|$ 以概率 1 趋于无穷——这是方差几何增长下随机游走的典型行为。

证明完毕。$\square$
\end{proof}

% ------------------------------------------------------------------
%  严格性暴击：RA.2
% ------------------------------------------------------------------
\begin{remark}[严格性暴击 / Honest Critique of Theorem RA.2]
\label{crit:ra2}
\honestcritique

定理 RA.2 的证明也有若干需要诚实审视的假设：

\begin{enumerate}[label=\textbf{暴击 \arabic*：}]
    \item \textbf{P\'olya 定理适用于 \iid 增量——但元审计噪声 $\xi_k$ 可能不是 \iid。} P\'olya 常返定理的核心假设是增量 $\{\xi_k\}$ 是独立同分布的。但在递归审计场景中，审计者可能从历史中学习：$\cA_{k+1}$ 在审视 $\cA_k$ 的输出时可能有意调整其估计策略，导致 $\xi_k$ 和 $\xi_{k+1}$ 相关。若 $\xi_k$ 具有负自相关（审计者矫枉过正），可能抑制随机游走的扩散，甚至使 $\rho=1$ 情形趋于收敛。若具有正自相关，发散可能更快。因此，严格依赖 \iid 假设可能高估或低估了实际发散风险。

    \item \textbf{$\rho=1$ 时 $f$ 为平凡的假设。} 实际中 $f$ 通常包含一定的收缩效应（即使 $\rho_f<1$ 可能不成立，$f$ 仍可能不是严格的恒等映射）。$f$ 的非平凡动力学与噪声的相互作用可能在 $\rho=1$ 时产生更丰富的相变行为，而非纯粹的随机游走。

    \item \textbf{中心极限定理的收敛速度。} 在论证 $P(|\hat{\eta}_k-\eta^*|<\varepsilon)\to 0$ 时，我们使用了 CLT 的近似。这在大 $k$ 下成立，但收敛速度由 Berry--Esseen 定理保证为 $O(1/\sqrt{k})$——这意味着尾部概率的衰减可能较慢，与定理声明的 ``趋于 0'' 不矛盾但提供了量化信息。

    \item \textbf{蒙特卡洛验证的可操作性。} 在实践中，$\rho=1$ 的随机游走行为可以通过模拟验证：若审计链中相邻层级的置信区间宽度大致相等（$\epsilon_k\approx\epsilon_{k+1}$），则系统处于临界点。但临界点的精确识别需要大量审计层级的观测，这在计算成本上可能是 prohibitive 的。
\end{enumerate}

综上，定理 RA.2 在 \iid 增量和纯随机游走的假设下是\textbf{严格}的（\rigorous）。但在实际审计系统中，$\xi_k$ 的非 \iid 性质可能改变结论，需逐例验证（\conditionallyrigorous）。
\end{remark}

% ------------------------------------------------------------------
\subsection{不动点信息论恒等式（Conjecture RA.1）}
% ------------------------------------------------------------------

\begin{conjecture}[不动点—审计热寂等价性 / Fixed-Point--Audit Heat Death Equivalence]
\label{conj:fixed-point}
递归审计不动点 $\eta^*$ 与 AE-定理中的审计热寂熵 $H_{\min}=H(\varepsilon\mid S)$ 相同：
\begin{equation}\label{eq:fixed-point-identity}
    \eta^* \;=\; H_{\min} \;=\; H(\varepsilon\mid S).
\end{equation}
若此猜想成立，则意味着：\textbf{无限递归审计的收敛点恰好是定理 3 的不可约无知（irreducible ignorance）。审计之剑指向自身时，最终刺中了自己锻造来刺穿的盾牌。}
\end{conjecture}

\medskip\noindent
\textbf{支持证据（Supporting Evidence）。}

尽管一般性证明尚未建立，以下证据为猜想提供方向性支持：

\begin{enumerate}[label=(\roman*)]

    \item \textbf{线性情形的显式匹配。}
    考虑线性审计算子 $f(\eta)=a\eta+b$，其中 $0\leq a<1$，$b\in[0,1-a]$ 以保持 $f:[0,1]\to[0,1]$。不动点方程为 $\eta^* = a\eta^*+b$，解得：
    \[
    \eta^* = \frac{b}{1-a}.
    \]
    另一方面，在简单伯努利噪声模型中，给定充分统计量 $S$ 下的条件熵 $H(\varepsilon\mid S)$ 也可写为类似形式。具体地，若 SCX 的 Cercis 评分系统将噪声过程建模为 $\varepsilon\sim\mathrm{Bernoulli}(\theta)$，且 $S$ 是共轭先验 Beta$(\alpha,\beta)$ 下的后验均值，则最小熵屏障为：
    \[
    H_{\min} = -\hat{\theta}\log\hat{\theta} - (1-\hat{\theta})\log(1-\hat{\theta}),
    \]
    其中 $\hat{\theta}=(\alpha+S)/(\alpha+\beta+n)$。存在参数化 $(a,b)$ 使得该熵值与 $b/(1-a)$ 匹配。

    \item \textbf{两者都是信息论算子的不动点。}
    $\eta^*$ 是递归审计算子 $f$ 的不动点：$\eta^* = f(\eta^*)$。
    $H_{\min}=H(\varepsilon\mid S)$ 是条件熵算子的不动点：在 SCX 框架中，信息获取过程 $\cI_A^{(\infty)}$ 作用后，$H(\varepsilon\mid S)$ 是经过无穷次审计信息注入后不可再减少的剩余熵。因此两者都刻画了 ``审计过程能到达的最远点''。

    \item \textbf{维度/量纲一致性。}
    在 SCX 框架中，$\eta^*$ 是极限噪声率（无量纲），$H_{\min}$ 是极限熵（以 nats 或 bits 为单位）。两者在数值上通过标准噪声率-熵转换公式 $\eta^* = h^{-1}(H_{\min})$ 关联，其中 $h(p)=-p\log p-(1-p)\log(1-p)$ 是二元熵函数。猜想实质上断言了这一转换是恒等映射——即不动点处的噪声率恰好等于其自身熵值的信息内容。

    \item \textbf{极限行为的唯一性。}
    由定理 RA.1，在 $\rho<1$ 条件下，$\eta^*$ 是递归审计序列的唯一极限。由 AE-定理，$H_{\min}$ 是 SCX 框架中任何审计过程的终极信息屏障。两者在各自体系内都是唯一的、吸引子的不动点——这为它们的同一性提供了拓扑学支持。

\end{enumerate}

\begin{remark}[猜想严格性标注 / Status]
\openquest

Conjecture \ref{conj:fixed-point} 目前是\textbf{开放猜想}。支持证据是实质性的但不完备——一般性证明需要建立 $f$ 作为 SCX 可观测量的 $\sigma$-代数上的审计-信息投影算子（audit-information projection operator）的严格表示定理。具体而言，需要证明：
\[
f(\eta) = \E[\varepsilon \mid \cI_A,\ \hat{\eta}=\eta],
\]
并进而证明其不动点与 $H(\varepsilon\mid S)$ 的条件极值问题等价。
\end{remark}

% ------------------------------------------------------------------
%  严格性暴击：Conjecture RA.1
% ------------------------------------------------------------------
\begin{remark}[诚实暴击 / Honest Critique of Conjecture RA.1]
\label{crit:conj}
\honestcritique

\begin{enumerate}[label=\textbf{暴击 \arabic*：}]

    \item \textbf{线性情形匹配不具一般性。} 证据 (i) 仅在线性算子下建立了匹配关系。SCX 框架中的审计算子 $f$ 可能是高度非线性的（如神经网络审计器），此时 $\eta^*=b/(1-a)$ 的显式形式不复存在，与 $H_{\min}$ 的参数匹配需要完全不同的论证。

    \item \textbf{信息论算子不等价于审计算子。} 证据 (ii) 指出两者都是不动点，但信息论算子 $\cI_A^{(\infty)}$ 和审计算子 $f$ 作用在完全不同的空间上（前者作用在熵值上，后者作用在噪声率上）。两者不动点的一致需要深层函子性（functoriality）证明。

    \item \textbf{二元熵函数 $h$ 的可逆性。} 证据 (iii) 的 $h^{-1}$ 映射在 $[0,1]$ 上不是单射——$h(p)=h(1-p)$，故 $H_{\min}$ 可能对应两个 $\eta^*$ 值。这需要额外的对称性破缺机制来唯一确定 $\eta^*$。

    \item \textbf{$\eta^*$ 的定义域依赖性。} 递归审计不动点 $\eta^*$ 依赖于初始审计器 $\cA_0$ 的选择（通过不动点方程的解），而 $H_{\min}$ 纯粹由数据分布和 SCX 评分定义。若 $\eta^*$ 依赖于初始条件（即 $f$ 有多个不动点），则恒等式只能在特定吸引域内成立。

\end{enumerate}
\end{remark}

% ------------------------------------------------------------------
\subsection{G\"odel 类比（开放问题）}
% ------------------------------------------------------------------

\begin{openproblem}[SCX G\"odel Correspondence / SCX G\"odel 对应]
\label{prob:godel}
是否存在以下形式对应关系：
\begin{equation}\label{eq:godel-correspondence}
    \text{``RA-定理中的 $\rho\geq1$''}
    \;\longleftrightarrow\;
    \text{``PA 中存在不可判定命题''},
\end{equation}
使得递归审计发散蕴含存在\textbf{绝对不可审计的审计命题（absolutely unauditable audit proposition）}——即关于审计质量的某个陈述，其真值无法通过 SCX 审计之剑的任何有限次迭代来确定？  \openquest
\end{openproblem}

\begin{remark}[G\"odel 类比的局限性 / Limitations of the G\"odel Analogy]
\label{rem:godel-limits}
\honestcritique

G\"odel 类比在哲学层面具有强烈吸引力，但需要诚实审视其形式局限：

\begin{enumerate}[label=\textbf{局限 \arabic*：}]

    \item \textbf{范畴差异。} G\"odel 不完备定理内在于\textbf{形式系统}的可证性概念：在一个包含皮亚诺算术的形式系统中，存在既不可证明也不可否证的语句。而 RA-定理涉及的是\textbf{随机过程}的收敛性质。将 ``不可证明性'' 类比为 ``非收敛性'' 是跨范畴的隐喻，而非严格的数学同构。

    \item \textbf{自指机制的差异。} G\"odel 的自指是通过哥德尔编号（公式的算术化）实现的元数学构造。RA-定理的自指是通过审计递归（审计者审计审计者...）实现的随机过程。前者是语法性的（syntactic），后者是统计性的（statistical）——两者的自指机制不同。

    \item \textbf{$\rho<1$ 对应什么？} 如果 $\rho\geq1$ 对应于 G\"odel 的不完备性，那么 $\rho<1$ 的收敛情形应对应于某个 ``可完备的形式系统''——但在 G\"odel 定理中，所有足够强大的形式系统都是不完备的。这个类比在 $\rho<1$ 情形下无法自然延伸。

    \item \textbf{绝对不可审计命题的存在性证明策略。} 若要在 SCX 框架内构造绝对不可审计命题，需要类比 G\"odel 的对角化论证：构造一个自称 ``我无法被审计到精度 $\varepsilon$ 以内'' 的审计结论。但 SCX 框架的概率性质使得这样的对角化论证面临自指悖论的连续化困难——审计结论是实数值而非真值。

\end{enumerate}

这些局限性不意味着 G\"odel 类比没有价值——它提出了深刻的问题：是否存在审计版本的 ``不可判定性''？但该问题的解决需要更多概念工作，可能涉及概率逻辑（probabilistic logic）或审计演算（audit calculus）的基础构建。 \openquest
\end{remark}

% ============================================================
\section{审计自指相图与讨论}
% ============================================================

\subsection{审计自指相图}

\begin{center}
\begin{tabular}{@{}llll@{}}
\toprule
$\rho$ 范围 & 相 & 审计之剑状态 & 保证程度 \\
\midrule
$\rho<1/2$    & 强收敛相    & 安全（指数收敛）      & \rigorous \\
$1/2\leq\rho<1$ & 弱收敛相 & 安全（多项式收敛）    & \rigorous \\
$\rho=1$      & 临界点      & 边缘（随机游走）      & \rigorous \\
$\rho>1$      & 发散相      & 断裂                  & \rigorous \\
\bottomrule
\end{tabular}
\end{center}

\subsection{严格性总结}

\begin{itemize}
    \item \textbf{Theorem RA.1 (递归收敛)：} \rigorous 在假设 \ref{as:contraction}--\ref{as:martingale-diff} 下完全严格。指数速率部分在 $\rho\leq1/2$ 条件下严格，但精确常数优化是开放问题。
    \item \textbf{Theorem RA.2 (无限递归)：} \rigorous 在 \iid 噪声和纯随机游走假设下严格。非 \iid 情形需要逐例验证。
    \item \textbf{Conjecture RA.1 (不动点恒等式)：} \openquest 支持证据充分但不完备。
    \item \textbf{G\"odel 类比：} \openquest 概念上的开放性远超数学上的形式化程度。
\end{itemize}

\subsection{工程含义}

条件 $\rho<1$ 要求每层元审计者严格优于其所审计的层级。这可通过以下方式实现：
\begin{enumerate}[label=(\roman*)]
    \item \textbf{人类升级：} $\cA_0$ = AI 审计者（大 $\sigma$），$\cA_1$ = 人类辅助审计者（小 $\sigma$），$\cA_2$ = 专家委员会（更小 $\sigma$）——自然达到 $\rho<1$；
    \item \textbf{数据增加：} 每层元审计者访问严格多于下层的额外数据；
    \item \textbf{更强模型：} 每层使用严格更强大的模型架构。
\end{enumerate}
核心洞见：\textbf{扁平审计层级（$\rho\approx1$）无法收敛。} SCX 部署需要具有严格递增保真度的\textbf{深层}审计层级。

\subsection{与其他定理的联系}

\begin{itemize}
    \item \textbf{TS-定理：} TS 的 Sprint 自审极限（漂移率上限）是 RA-定理 $\rho$ 条件的时间版本；漂移率 $v_{\mathrm{drift}}$ 对应递归序列中的每步 $\sigma_k^2$。
    \item \textbf{AE-定理：} 猜想 \ref{conj:fixed-point} 将 $\eta^*$ 与 $H_{\min}$ 统一——递归收敛点与热力学热寂点。
    \item \textbf{HC-定理：} 人类干预对应在审计层级中注入 $\rho\ll1$ 的层级，打破等保真度死锁。
    \item \textbf{FA-定理：} 联邦审计层级要求 $\prod_k \rho_k < 1$ 以跨联邦层级全局收敛。
\end{itemize}

% ============================================================
\section{Conclusion}
% ============================================================

RA-Theorem 将审计之剑置于自我反射之下。答案是微妙的：
\begin{enumerate}[label=(\arabic*)]
    \item 当元审计者严格更高保真度时（$\rho<1$），递归收敛——剑成立。  \rigorous
    \item 当元审计者等保真度或更低保真度时（$\rho\geq1$），递归发散——剑断裂。  \rigorous
    \item 收敛点被猜想为定理 3 的屏障——$H_{\min}$，即不可约无知。  \openquest
    \item G\"odel 对应——发散递归是否蕴含绝对不可审计命题——是基础性的开放问题。  \openquest
\end{enumerate}

实践指令清晰明确：\textbf{永远不要部署审计者与元审计者等保真度的自审计系统。} 审计之剑需要层级，而非循环。

% ============================================================
\bibliographystyle{plainnat}
\begin{thebibliography}{30}

\bibitem{scx2024cercis}
SCX Working Group.
\newblock ``The SCX Axiomatic Framework: Cercis, Situs, and Tiresias
Operators,''
\newblock \emph{Technical Report}, 2024.

\bibitem{godel1931formal}
K.~G\"odel.
\newblock ``\"Uber formal unentscheidbare S\"atze der Principia Mathematica
und verwandter Systeme I,''
\newblock \emph{Monatshefte f\"ur Mathematik und Physik}, 38:173--198, 1931.

\bibitem{turing1936computable}
A.~M.~Turing.
\newblock ``On computable numbers, with an application to the
Entscheidungsproblem,''
\newblock \emph{Proceedings of the London Mathematical Society},
s2-42(1):230--265, 1937.

\bibitem{banach1922operations}
S.~Banach.
\newblock ``Sur les op\'erations dans les ensembles abstraits et leur
application aux \'equations int\'egrales,''
\newblock \emph{Fundamenta Mathematicae}, 3:133--181, 1922.

\bibitem{doob1953stochastic}
J.~L.~Doob.
\newblock \emph{Stochastic Processes}.
\newblock Wiley, 1953.

\bibitem{polya1921random}
G.~P\'olya.
\newblock ``\"Uber eine Aufgabe der Wahrscheinlichkeitsrechnung betreffend die
Irrfahrt im Stra\ss ennetz,''
\newblock \emph{Mathematische Annalen}, 84:149--160, 1921.

\bibitem{meyn2012markov}
S.~P.~Meyn and R.~L.~Tweedie.
\newblock \emph{Markov Chains and Stochastic Stability}, 2nd ed.
\newblock Cambridge University Press, 2012.

\bibitem{hajek1965martingale}
J.~H\'ajek and A.~R\'enyi.
\newblock ``Generalization of an inequality of Kolmogorov,''
\newblock \emph{Acta Mathematica Academiae Scientiarum Hungaricae},
6(3--4):281--283, 1955.

\bibitem{smullyan1992godel}
R.~M.~Smullyan.
\newblock \emph{G\"odel's Incompleteness Theorems}.
\newblock Oxford University Press, 1992.

\bibitem{hofstadter1979godel}
D.~R.~Hofstadter.
\newblock \emph{G\"odel, Escher, Bach: An Eternal Golden Braid}.
\newblock Basic Books, 1979.

\bibitem{cover1999elements}
T.~M.~Cover and J.~A.~Thomas.
\newblock \emph{Elements of Information Theory}, 2nd ed.
\newblock Wiley, 2006.

\bibitem{villani2008optimal}
C.~Villani.
\newblock \emph{Optimal Transport: Old and New}.
\newblock Springer, 2008.

\end{thebibliography}

\end{document}
